\scanpage{101}

Je préfère faire \unsure{le calcul|les calculs} du côté de $A_1,A_2$. On peut expliciter la construction de $A$ ainsi

\diagram{\begin{tikzcd}[ampersand replacement=\&,column sep=small]
0\arrow[r] \& \overbrace{A_1\deleted{\unsure{\oplus}}\amalg_{\underline{\mathcal O}}A_2}^{=A_1\otimes_0 A_2}\arrow[r,dash]\arrow[d,"\tau"'] \& A_1\otimes A_2\arrow[r]\arrow[d,dash] \& L_1^\vee\otimes L_2^\vee\arrow[r]\arrow[d,equal] \& 0\\
0\arrow[r] \& \underline{\mathcal O}\arrow[r] \& A\arrow[r] \& L_1^\vee\otimes L_2^\vee\arrow[r] \& 0
\end{tikzcd}}

\margin{\unsure{Trace réduite|Tr évidente}, qui est $\mathrm{Tr}_{A_1},\mathrm{Tr}_{A_2}$ sur les deux composantes}

$A$ \inserted{$=A_1*A_2$} est l'ext. de $L_1^\vee\otimes L_2^\vee$ par $\underline{\mathcal O}$, déduite de l'extension $A_1\otimes A_2$ de $L_1^\vee\otimes L_{2}^{\unsure{\vee|\prime}}$ par $A_1\otimes_0 A_2$, par l'hom $\tau$

\[
A_1\amalg_{\underline{\mathcal O}}A_2\longrightarrow\mathcal O.
\]

Pour donner l'hom

\[
A\longrightarrow A_1\otimes A_2
\]

c'est se donner séparément

\[
\underline{\mathcal O}\longrightarrow A_1\otimes A_2
\]

(on prendra l'inclusion canonique $\lambda\longmapsto\lambda.e_{A_1}\otimes e_{A_2}$

et

\[
A_1\otimes A_2\xrightarrow[x_1\otimes x_2\mapsto i(x_1*x_2)]{\unsure{\leftrightharpoons|\leftrightarrows|\varphi|\natural}}A_1\otimes A_{\unsure{2}}
\]

tel que les composés avec

\diagram{\begin{tikzcd}[ampersand replacement=\&,column sep=small]
\&\mathcal O\\
\deleted{\unsure{E_1}}A_1\otimes_0\deleted{\unsure{E_2}}A_2=A_1\oplus_{\mathcal O}A_2\arrow[ur]\arrow[dr]\&\\
\& A_1\otimes A_2
\end{tikzcd}}

soient égaux. On prendra

% archive sheet number: 100

